\documentclass[14pt,a4paper,oneside]{extbook}
\usepackage[utf8]{vietnam}
\usepackage{amsmath,amssymb,mathrsfs,fancyhdr,enumerate,multirow,makecell}
\usepackage{tkz-euclide,tikz-3dplot,tikz,tkz-tab}
\usetikzlibrary{arrows.meta,calc,intersections,angles,quotes,backgrounds,shapes.geometric,patterns,through,shadows,positioning}
\usepackage{yhmath}
\usepackage[top=0.7cm, bottom=0.7cm, left=1.0cm, right=1.0cm, headheight=16pt, headsep=3mm, footskip=5mm]{geometry}
\usepackage[hidelinks,unicode]{hyperref}
\usepackage{qrcode}
\usepackage[many]{tcolorbox}
\usepackage{fontawesome}

\pagestyle{fancy}
\fancyhf{}
\fancyhead[L]{\footnotesize\bfseries\color{teal!85!black}\faBook~SỔ TAY KIẾN THỨC TOÁN 5}
\fancyhead[R]{\footnotesize\bfseries\color{blue!85!black}Lớp Toán Thầy Lĩnh -- loptoanthaylinh.vn}
\fancyfoot[C]{\footnotesize\bfseries\color{gray!80!black}-- Sổ tay kiến thức: Bài 66. Thực hành trải nghiệm thống kê --}
\renewcommand{\headrulewidth}{0.4pt}

\begin{document}

% --- BANNER TIÊU ĐỀ BÀI HỌC ---
\noindent
\begin{tcolorbox}[
  enhanced,
  colback=teal!8!white,
  colframe=teal!80!black,
  arc=3mm,
  boxrule=1.1pt,
  width=\linewidth,
  halign=center,
  top=2.2mm, bottom=2.2mm,
  left=3mm, right=3mm,
  drop fuzzy shadow
]
  {\large\bfseries\color{teal!95!black}\faBookmark\ BÀI 66. THỰC HÀNH TRẢI NGHIỆM THỐNG KÊ}\\[2pt]
  {\scriptsize\color{black!75}\itshape Các bước tiến hành một cuộc điều tra thống kê thực tế, xử lý và báo cáo kết quả trên 1 trang A4}
\end{tcolorbox}

\vspace{3.2mm}
% --- KHỐI 1: QUY TRÌNH THỰC HIỆN ĐIỀU TRA ---
\noindent
\begin{minipage}[t]{0.56\textwidth}
  \begin{tcolorbox}[
    enhanced,
    title={\footnotesize\bfseries\color{white}\faTasks\ 1. Bốn giai đoạn điều tra thống kê},
    colback=blue!2!white,
    colframe=blue!75!black,
    arc=2.5mm, boxrule=0.8pt, top=1.8mm, bottom=1.8mm, left=2.2mm, right=2.2mm
  ]
    \footnotesize
    $\bullet$ \textbf{Giai đoạn 1 (Chuẩn bị):} Xác định đề tài, đối tượng và thiết kế phiếu hỏi.\\
    $\bullet$ \textbf{Giai đoạn 2 (Thu thập):} Phát phiếu, phỏng vấn và kiểm đếm gạch sổ.\\
    $\bullet$ \textbf{Giai đoạn 3 (Xử lý):} Lập bảng tổng hợp, tính tỉ lệ phần trăm.\\
    $\bullet$ \textbf{Giai đoạn 4 (Báo cáo):} Vẽ biểu đồ và đưa ra kết luận, nhận xét.
  \end{tcolorbox}
\end{minipage}%
\hfill
\begin{minipage}[t]{0.42\textwidth}
  \begin{tcolorbox}[
    enhanced,
    title={\footnotesize\bfseries\color{white}\faBullhorn\ 2. Kĩ năng nhận xét dữ liệu},
    colback=teal!2!white,
    colframe=teal!75!black,
    arc=2.5mm, boxrule=0.8pt, top=1.8mm, bottom=1.8mm, left=2mm, right=2mm
  ]
    \footnotesize
    $\bullet$ \textbf{Đánh giá xu hướng:} Xu thế nào được yêu thích/lựa chọn nhiều nhất?\\
    $\bullet$ \textbf{Đề xuất giải pháp:} Dựa vào số liệu thực tế để đưa ra gợi ý hợp lý trong học tập và đời sống.
  \end{tcolorbox}
\end{minipage}

\vspace{3mm}

% --- KHỐI 3: SƠ ĐỒ TƯ DUY ---
\noindent
\begin{tcolorbox}[
  enhanced,
  title={\footnotesize\bfseries\color{white}\faSitemap\ Sơ đồ tư duy: Quy trình điều tra thống kê trải nghiệm},
  colback=yellow!2!white,
  colframe=teal!85!black,
  arc=3mm, boxrule=0.9pt, top=1.8mm, bottom=1.8mm, left=1.5mm, right=1.5mm,
  halign=center
]
  \begin{tikzpicture}[
    scale=0.92,
    every node/.style={transform shape},
    s1/.style={rectangle, rounded corners=4pt, fill=blue!15, draw=blue!70!black, text=blue!90!black, font=\bfseries\scriptsize, inner sep=4pt, align=center},
    s2/.style={rectangle, rounded corners=4pt, fill=teal!15, draw=teal!70!black, text=teal!90!black, font=\bfseries\scriptsize, inner sep=4pt, align=center},
    s3/.style={rectangle, rounded corners=4pt, fill=orange!15, draw=orange!70!black, text=orange!90!black, font=\bfseries\scriptsize, inner sep=4pt, align=center},
    s4/.style={rectangle, rounded corners=4pt, fill=green!15, draw=green!70!black, text=green!90!black, font=\bfseries\scriptsize, inner sep=4pt, align=center}
  ]
    \node[s1] (n1) at (-6.3, 0) {1. Thiết kế\\Phiếu điều tra};
    \node[s2] (n2) at (-2.1, 0) {2. Thu thập \&\\Kiểm đếm gạch sổ};
    \node[s3] (n3) at (2.1, 0) {3. Lập bảng số liệu\\\& Tính tỉ số \%};
    \node[s4] (n4) at (6.3, 0) {4. Báo cáo \&\\Đưa ra đề xuất};

    \draw[teal!85!black, thick, -{Stealth[length=2.5mm]}] (n1) -- (n2);
    \draw[teal!85!black, thick, -{Stealth[length=2.5mm]}] (n2) -- (n3);
    \draw[teal!85!black, thick, -{Stealth[length=2.5mm]}] (n3) -- (n4);
  \end{tikzpicture}
\end{tcolorbox}

\vspace{3mm}

% --- KHỐI 4: KỸ THUẬT & BẪY TRẮC NGHIỆM ---
\noindent
\begin{minipage}[t]{0.54\textwidth}
  \begin{tcolorbox}[
    enhanced,
    title={\footnotesize\bfseries\color{white}\faCheckCircle\ Checklist chất lượng số liệu},
    colback=green!3!white,
    colframe=green!60!black,
    arc=2.5mm, boxrule=0.8pt, top=1.8mm, bottom=1.8mm, left=2.2mm, right=2.2mm
  ]
    \footnotesize
    $\square$ Phiếu có bị tích trùng $2$ phương án không?\\
    $\square$ Có phiếu nào bị để trống câu trả lời không?\\
    $\square$ Tổng số phiếu thu về có khớp với sĩ số lớp không?\\
    $\square$ Các số liệu thành phần cộng lại có bằng $100\%$ không?
  \end{tcolorbox}
\end{minipage}%
\hfill
\begin{minipage}[t]{0.44\textwidth}
  \begin{tcolorbox}[
    enhanced,
    title={\footnotesize\bfseries\color{white}\faExclamationTriangle\ Mẹo tránh bẫy trắc nghiệm},
    colback=red!2!white,
    colframe=red!70!black,
    arc=2.5mm, boxrule=0.8pt, top=1.8mm, bottom=1.8mm, left=2.2mm, right=2.2mm
  ]
    \footnotesize
    $\bullet$ \textbf{Lỗi kết luận vội vã:} Không được đưa ra kết luận khi cỡ mẫu điều tra quá nhỏ (ví dụ chỉ hỏi $2$ bạn).\\
    $\bullet$ Luôn kiểm tra tính khách quan của câu hỏi điều tra.
  \end{tcolorbox}
\end{minipage}

\end{document}
